\chapter{How Duq Works}
\label{ch:how_it_works}

Duq operates through a high-precision DSP engine that maps MIDI input to volume modulation.

\section{Trigger Engine}
When Duq receives a MIDI Note On message, it identifies the assigned envelope for that specific MIDI note. The plugin supports multiple envelopes, each mapped to a unique MIDI trigger note.

\section{Envelope Execution}
Once triggered, the envelope begins its playback cycle. Duq supports two primary timing modes:
\begin{enumerate}
    \item \textbf{Frequency Mode (Hz):} The envelope duration is determined by a fixed frequency. In this mode, if the MIDI note is held, the envelope will loop. If the note is released, the envelope finishes its current cycle and stops.
    \item \textbf{Sync Mode (BPM):} The envelope duration is synchronized to the DAW's tempo and time signature (e.g., 1/4 note, 1/8 note).
\end{enumerate}

\section{Signal Path}
1. \textbf{Input:} The audio signal enters the plugin.
2. \textbf{Lookahead Buffer:} The signal is delayed by a user-defined lookahead time (0--100ms).
3. \textbf{Modulation:} The triggered envelope shapes the gain of the delayed signal.
4. \textbf{Mix:} The processed (wet) signal is blended with the original (dry) signal.
5. \textbf{Output:} The final blended signal is sent to the DAW.
